\documentclass[11pt, a4paper]{article}

\usepackage[utf8]{inputenc}
\usepackage[T1]{fontenc}
\usepackage[french]{babel}
\usepackage{amsmath, amssymb}
\usepackage{geometry}
\usepackage{graphicx}
\usepackage{tabularx}
\usepackage{booktabs}
\usepackage{xcolor}
\usepackage{listings}
\usepackage{tikz}
\usepackage{hyperref}

% Configuration des marges
\geometry{
    top=2cm,
    bottom=2cm,
    left=1.5cm,
    right=1.5cm
}

% Configuration pour les blocs de code VHDL
\lstset{
    basicstyle=\ttfamily\small,
    breaklines=true,
    frame=single,
    tabsize=4,
    showstringspaces=false,
    keywordstyle=\color{blue}\bfseries,
    commentstyle=\color{green!50!black},
    stringstyle=\color{red},
    numbers=left,
    numberstyle=\tiny\color{gray}
}

% Commande pour les placeholders d'images
\newcommand{\imageplaceholder}[2]{
    \begin{center}
        \fbox{\begin{minipage}{#1}\centering\vspace{1cm}\textit{#2}\vspace{1cm}\end{minipage}}
    \end{center}
}

% Commande pour afficher la correction avec explication
\newcommand{\corr}[1]{
    \vspace{0.3cm}
    \begin{center}
    \fcolorbox{blue}{blue!5}{
    \begin{minipage}{0.95\textwidth}
    \color{blue} \textbf{Correction \& Raisonnement :} \\
    #1
    \end{minipage}}
    \end{center}
    \vspace{0.3cm}
}

\begin{document}

% ==================== EN-TÊTE ====================
\noindent
\begin{tabularx}{\textwidth}{@{}X c r@{}}
\textbf{SN361} & \textbf{Session 1} & \textbf{Documents selon sujet} \\
\textbf{VHDL} & \textbf{Examen} & \textbf{Calculatrice selon sujet} \\
\end{tabularx}
\vspace{0.2cm}
\hrule
\vspace{0.4cm}

% ==================== PROBLÈME I ====================
\begin{center}
    \textbf{\Large Problème I - Conception et validation d'un « timer » - 90 pts}
\end{center}
\vspace{-0.2cm}
\hrule
\vspace{0.3cm}

\subsubsection*{Contexte}
\noindent L'objectif est de concevoir un timer générant périodiquement une impulsion d'une période d'horloge de large. Ce timer est contrôlé par une horloge interne de période $1\,\mu\text{s}$. Il a une largeur de 32 bits.

\subsubsection*{Lecture du Schéma : Liste des éléments et leur utilité}
\noindent Le schéma illustre l'architecture matérielle (chemin de données et contrôle) du timer. Voici comment le lire pas-à-pas :
\begin{itemize}
    \item \textbf{1 microsec oscillator (Oscillateur)} : Fournit le signal d'horloge de base (période de $1\,\mu\text{s}$) qui cadence tous les éléments synchrones du circuit.
    \item \textbf{Bloc "-1" (Décrémenteur combinatoire)} : Réalise une opération de soustraction de 1 sur l'entrée $M$. Cela est nécessaire car le compteur décrémente jusqu'à 0 inclus. Pour compter $M$ cycles, il faut démarrer le comptage à $M-1$.
    \item \textbf{32-bit register (Registre de mémorisation)} : Sauvegarde la valeur $M-1$ à chaque fois que la commande \texttt{load} est active (à '1'). Il sert de mémoire pour conserver la configuration du timer de manière autonome.
    \item \textbf{4-bit 2x1 Mux (Multiplexeur 2 vers 1)} : Aiguille la donnée qui sera chargée dans le compteur. Si \texttt{load} est à '1', il laisse passer la nouvelle valeur $M-1$. Si \texttt{load} est à '0', il laisse passer la valeur sauvegardée dans le registre.
    \item \textbf{Porte logique OU (OR gate)} : Génère le signal de chargement (\texttt{ld}) du compteur (si \texttt{load = 1} ou \texttt{tc = 1}).
    \item \textbf{32-bit down-counter (Décompteur 32 bits)} : C'est le cœur du timer. Activé par \texttt{enable}, il décrémente sa valeur à chaque front d'horloge. Lorsqu'il atteint 0, sa sortie \texttt{tc} passe à 1.
\end{itemize}

\vspace{0.3cm}
\noindent \textbf{Question 1} : Quelle valeur doit être programmée (c'est-à-dire envoyée vers M) pour générer une impulsion toutes les 300 ms ?

\corr{
\textbf{Étape 1 : Analyse des périodes.}\\
L'horloge a une période de $1\,\mu\text{s}$. On souhaite une période globale de $300\text{ ms}$, ce qui équivaut à $300\,000\,\mu\text{s}$.

\textbf{Étape 2 : Calcul du nombre de cycles.}\\
Il faut donc compter exactement $300\,000$ cycles d'horloge entre chaque impulsion.

\textbf{Étape 3 : Logique de comptage.}\\
Le circuit décrémente de $M-1$ jusqu'à $0$ (soit un total de $M$ états ou $M$ cycles). Le bloc "-1" matériel s'occupe déjà de soustraire 1 à la valeur de consigne.

\textbf{Étape 4 : Déduction de M.}\\
La valeur à programmer sur l'entrée $M$ correspond donc directement au nombre de cycles souhaités :
$$ \mathbf{M = 300\,000} $$
}

\noindent \textbf{Question 2} : Faire un chronogramme pour deux impulsions avec $M=3$ puis deux impulsions avec $M=2$.

\corr{
\textbf{Étape 1 : Initialisation (Cycle 0).}\\
\texttt{load} = 1, $M = 3$. Le bloc "-1" calcule $M-1 = 2$. Le registre R32 charge 2. Le Mux32 laisse passer 2. Le compteur C charge 2 car la porte OU donne 1.

\textbf{Étape 2 : Première séquence d'impulsions (M=3).}\\
\begin{itemize}
    \item \textbf{Cycle 1 :} \texttt{load} = 0, \texttt{enable} = 1. C = 2. \texttt{tc} (Q) = 0. Au front, C $\rightarrow$ 1.
    \item \textbf{Cycle 2 :} C = 1, Q = 0. Au front, C $\rightarrow$ 0.
    \item \textbf{Cycle 3 (Impulsion) :} C = 0. Donc Q = 1 (\texttt{tc}=1). La porte OU vaut 1. Le Mux32 laisse passer R32 (2). Au front, C recharge 2.
    \item \textbf{Cycles 4, 5, 6 :} Identique aux cycles 1 à 3. On obtient une 2ème impulsion.
\end{itemize}

\textbf{Étape 3 : Seconde séquence d'impulsions (M=2).}\\
\begin{itemize}
    \item \textbf{Cycle 7 (Changement) :} \texttt{load} = 1, $M = 2$. R32 et C chargent 1 ($2-1$).
    \item \textbf{Cycle 8 (Impulsion) :} \texttt{load} = 0, \texttt{enable} = 1. C passe de 1 à 0 au cycle 9. La période est maintenant de 2 cycles.
\end{itemize}
}

\noindent \textbf{Question 3} : Est-ce que les périodes des impulsions générées correspondent aux valeurs de M successives (3 puis 2) ? Si ce n'est pas le cas, expliquez.

\corr{
\textbf{Étape 1 : Vérification pour M=3.}\\
Pour $M=3$, le compteur prend les valeurs $2 \rightarrow 1 \rightarrow 0$. Il faut donc 3 cycles d'horloge complets pour revenir à l'état initial. La période correspond.

\textbf{Étape 2 : Vérification pour M=2.}\\
Pour $M=2$, le compteur prend les valeurs $1 \rightarrow 0$. Il faut donc 2 cycles d'horloge. La période correspond.

\textbf{Étape 3 : Conclusion et rôle du bloc soustracteur.}\\
\textbf{Oui}, les périodes correspondent parfaitement. Le bloc "-1" placé en amont compense le fait que le compteur passe par l'état "0". Il garantit que la période mesurée en cycles d'horloge soit exactement égale à la valeur entière $M$ demandée.
}

\noindent \textbf{Question 4} : Donnez une description en VHDL de l'entité du timer.

\corr{
\textbf{Étape 1 : Identification des ports et bibliothèques.}\\
Le schéma montre les entrées d'horloge, de contrôle (\texttt{load}, \texttt{enable}), le bus de données $M$ sur 32 bits et la sortie scalaire $Q$. On inclut `numeric_std` en prévision de l'architecture.
\begin{lstlisting}[language=VHDL]
library ieee;
use ieee.std_logic_1164.all;
use ieee.numeric_std.all;

entity timer is
    port (
        clk    : in  std_logic;
        load   : in  std_logic;
        enable : in  std_logic;
        M      : in  std_logic_vector(31 downto 0);
        Q      : out std_logic
    );
end entity timer;
\end{lstlisting}
}

\noindent \textbf{Question 5} : Décrire en VHDL RTL l'architecture \texttt{rtl} de ce circuit.

\corr{
\textbf{Étape 1 : Déclaration des signaux internes.}\\
On déclare les signaux de type `unsigned` pour les bus de données (M32, R32, Mux32, C) pour permettre l'arithmétique.

\textbf{Étape 2 : Implémentation combinatoire et processus synchrones.}\\
La logique combinatoire gère le soustracteur, le multiplexeur, la porte OU et le comparateur à zéro. Deux processus distincts gèrent les éléments mémorisants (Registre R32 et Compteur C).
\begin{lstlisting}[language=VHDL]
architecture rtl of timer is
    signal M32_uns, R32_uns, Mux32_uns, C_uns : unsigned(31 downto 0);
    signal OR1, Q_int : std_logic;
begin
    -- Logique Combinatoire
    M32_uns <= unsigned(M) - 1;
    Mux32_uns <= M32_uns when load = '1' else R32_uns;
    
    Q_int <= '1' when C_uns = to_unsigned(0, 32) else '0';
    OR1 <= load or Q_int;
    Q <= Q_int;

    -- Registre de Memorisation (R32)
    process(clk) begin
        if rising_edge(clk) then
            if load = '1' then
                R32_uns <= M32_uns;
            end if;
        end if;
    end process;

    -- Decompteur 32 bits (C)
    process(clk) begin
        if rising_edge(clk) then
            if OR1 = '1' then
                C_uns <= Mux32_uns;     -- Chargement synchrone
            elsif enable = '1' then
                C_uns <= C_uns - 1;     -- Decomptage synchrone
            end if;
        end if;
    end process;
end architecture rtl;
\end{lstlisting}
}

\noindent \textbf{Question 6} : Commentez la fonction du code Verilog fourni. \textit{(Note: Fait référence à un bout de code Verilog issu du sujet original).}

\corr{
\textbf{Étape 1 : Traduction du comportement.}\\
Le code Verilog décrit un \textbf{multiplexeur 2 vers 1} modélisé de manière combinatoire (sensibilité \texttt{always @(*)} ou \texttt{always @(load or M32 or R32)}).

\textbf{Étape 2 : Conclusion.}\\
Il aiguille le signal : si \texttt{load} vaut 1, la sortie prend la valeur de l'entrée \texttt{M32}. Si \texttt{load} vaut 0, la sortie prend la valeur de l'entrée \texttt{R32}.
}

\noindent \textbf{Question 7} : Combien de bascules générera la synthèse de votre architecture précédente ? Justifiez.

\corr{
\textbf{Étape 1 : Analyse des processus synchrones.}\\
Le circuit contient deux processus réagissant au front montant de l'horloge (\texttt{rising\_edge(clk)}), qui vont donc synthétiser des bascules D.

\textbf{Étape 2 : Dénombrement.}\\
\begin{itemize}
    \item Le premier processus mémorise le vecteur \texttt{R32\_uns} (32 bits) $\rightarrow$ 32 bascules.
    \item Le second processus mémorise le vecteur \texttt{C\_uns} (32 bits) $\rightarrow$ 32 bascules.
\end{itemize}
Tous les autres signaux (Mux, bloc -1, porte OU) sont modélisés en combinatoire.
$$ \mathbf{Total = 64 \text{ bascules}} $$
}

\noindent \textbf{Question 8} : Donnez une description du testbench \texttt{tb\_timer} permettant de retrouver exactement le même chronogramme que celui de la question 2.

\corr{
\textbf{Étape 1 : Instanciation et horloge.}\\
On crée l'entité de test, on génère une horloge de $1\,\mu\text{s}$ et on instancie le \texttt{timer}.

\textbf{Étape 2 : Scénario (Stimuli).}\\
\begin{lstlisting}[language=VHDL]
library ieee;
use ieee.std_logic_1164.all;
use ieee.numeric_std.all;

entity tb_timer is
end entity tb_timer;

architecture sim of tb_timer is
    signal clk, load, enable, Q : std_logic := '0';
    signal M : std_logic_vector(31 downto 0) := (others => '0');
    constant CLK_PERIOD : time := 1 us;
begin
    UUT: entity work.timer port map(clk, load, enable, M, Q);

    clk_process : process begin
        clk <= '0'; wait for CLK_PERIOD / 2;
        clk <= '1'; wait for CLK_PERIOD / 2;
    end process;

    stim_proc : process begin
        -- Initialisation (Cycle 0)
        M <= std_logic_vector(to_unsigned(3, 32));
        load <= '1'; enable <= '0'; wait for CLK_PERIOD;
        
        -- Generation de 2 impulsions (M=3 -> 6 cycles)
        load <= '0'; enable <= '1'; wait for 6 * CLK_PERIOD;
        
        -- Changement (M=2)
        M <= std_logic_vector(to_unsigned(2, 32));
        load <= '1'; wait for CLK_PERIOD;
        
        -- Generation de 2 impulsions (M=2 -> 4 cycles)
        load <= '0'; wait for 4 * CLK_PERIOD;
        
        enable <= '0'; wait;
    end process;
end architecture sim;
\end{lstlisting}
}

\newpage

% ==================== PROBLÈME II ====================
\begin{center}
    \textbf{\Large Problème II - Contrôleur de porte de garage - 30 pts}
\end{center}
\vspace{-0.2cm}
\hrule
\vspace{0.3cm}

\subsubsection*{Contexte}
\noindent Une porte de garage est contrôlée par une machine à états. \texttt{ctr(1)} gère l'alimentation du moteur (1=ON, 0=OFF) et \texttt{ctr(0)} définit le sens (0=ouverture, 1=fermeture). Les capteurs \texttt{sen1} et \texttt{sen2} détectent respectivement l'ouverture totale et la fermeture totale. Le bouton \texttt{remt} donne les ordres.

\vspace{0.3cm}
\noindent \textbf{Question 2.1} : Donnez un diagramme de transitions de la machine à états de Moore avec \texttt{remt} sous forme d'impulsion d'une seule période.

\corr{
\textbf{Étape 1 : Définition des états et sorties.}\\
La sortie dépend uniquement de l'état (Machine de Moore). On définit les vecteurs \texttt{ctr} associés :
\begin{itemize}
    \item \textbf{FERME} (\texttt{ctr = "0X"}) : Moteur coupé.
    \item \textbf{OUVERTURE} (\texttt{ctr = "10"}) : Moteur ON, sens ouverture.
    \item \textbf{OUVERT} (\texttt{ctr = "0X"}) : Moteur coupé.
    \item \textbf{FERMETURE} (\texttt{ctr = "11"}) : Moteur ON, sens fermeture.
    \item \textbf{STOP\_OUVERTURE} (\texttt{ctr = "0X"}) : Arrêt d'urgence/pause.
    \item \textbf{STOP\_FERMETURE} (\texttt{ctr = "0X"}) : Arrêt d'urgence/pause.
\end{itemize}

\textbf{Étape 2 : Conditions de transition.}\\
\begin{itemize}
    \item Depuis \textbf{FERME} : Si \texttt{remt = 1} $\rightarrow$ \textbf{OUVERTURE}.
    \item Depuis \textbf{OUVERTURE} : Si \texttt{sen1 = 1} $\rightarrow$ \textbf{OUVERT}. Si \texttt{remt = 1} $\rightarrow$ \textbf{STOP\_OUVERTURE}.
    \item Depuis \textbf{OUVERT} : Si \texttt{remt = 1} $\rightarrow$ \textbf{FERMETURE}.
    \item Depuis \textbf{FERMETURE} : Si \texttt{sen2 = 1} $\rightarrow$ \textbf{FERME}. Si \texttt{remt = 1} $\rightarrow$ \textbf{STOP\_FERMETURE}.
    \item Depuis \textbf{STOP\_OUVERTURE} : Si \texttt{remt = 1} $\rightarrow$ \textbf{FERMETURE} (inversion de sens).
    \item Depuis \textbf{STOP\_FERMETURE} : Si \texttt{remt = 1} $\rightarrow$ \textbf{OUVERTURE} (inversion de sens).
\end{itemize}
}

\noindent \textbf{Question 2.2} : Donnez un diagramme de transitions de la machine à états de Moore si le signal \texttt{remt} est un signal continu tant qu'on appuie sur le bouton. Il faut détecter les passages à 0.

\corr{
\textbf{Étape 1 : Identification du problème (Bouton continu).}\\
Si \texttt{remt} reste à '1' sur plusieurs cycles d'horloge, la machine basculerait instantanément entre, par exemple, \textbf{OUVERTURE} et \textbf{STOP\_OUVERTURE} à chaque coup d'horloge.

\textbf{Étape 2 : Résolution avec des états d'attente (Wait states).}\\
Pour éviter le rebouclage indésirable, il faut attendre le relâchement du bouton (\texttt{remt = 0}) avant de valider l'action. On double les états déclenchés par \texttt{remt} :
\begin{itemize}
    \item \textbf{FERME} : Si \texttt{remt = 1} $\rightarrow$ \textbf{ATTENTE\_RELACHEMENT\_OUV} (\texttt{ctr = "0X"}).
    \item \textbf{ATTENTE\_RELACHEMENT\_OUV} : Si \texttt{remt = 0} $\rightarrow$ \textbf{OUVERTURE}. Si \texttt{remt = 1}, boucle sur lui-même.
    \item \textbf{OUVERTURE} : Si \texttt{remt = 1} $\rightarrow$ \textbf{ATTENTE\_RELACHEMENT\_STOP} (\texttt{ctr = "0X"}).
    \item \textbf{ATTENTE\_RELACHEMENT\_STOP} : Si \texttt{remt = 0} $\rightarrow$ \textbf{STOP\_OUVERTURE}.
\end{itemize}
Cette logique s'applique symétriquement pour la fermeture.
}

\vfill
% ==================== PIED DE PAGE ====================
\hrule
\vspace{0.1cm}
\noindent
Équipe Pédagogique \hfill Esisar \hfill Page 1/1

\end{document}